\subsection{Ellipticity and minimal Laplace-type operator}
\label{subsec:ellipticity}

Let $(M,g)$ be a smooth four-dimensional Riemannian manifold.
We consider a second-order differential operator acting on scalar functions.

The analysis relies crucially on ellipticity.
For a differential operator of order two, ellipticity means that its principal symbol is invertible away from the zero section.

We choose the minimal Laplace-type operator
\begin{equation}
  A_g = -\nabla_g^2 ,
\end{equation}
where $\nabla_g$ is the Levi-Civita connection.
Its principal symbol is
\begin{equation}
  \sigma_2(A_g)(x,k) = g^{\mu\nu}(x) k_\mu k_\nu .
\end{equation}
On a Riemannian manifold, $g^{\mu\nu}k_\mu k_\nu>0$ for all nonzero $k$.
Therefore $A_g$ is elliptic.

Ellipticity ensures:

\begin{itemize}
  \item existence of the resolvent $(A_g+\lambda)^{-1}$ for $\lambda>0$,
  \item discreteness of the spectrum on compact manifolds,
  \item local regularity of the Green kernel,
  \item validity of the short-time heat-kernel expansion,
  \item meromorphic continuation of the spectral zeta function.
\end{itemize}

More general scalar Laplace-type operators may include a curvature coupling,
\begin{equation}
  A_g(\xi) = -\nabla_g^2 + \xi R .
\end{equation}
In this work we set $\xi=0$.
This is the minimal choice in which the geometry enters exclusively through the principal symbol.
The conformal value $\xi=1/6$, for which the Einstein--Hilbert counterterm vanishes in four dimensions, is discussed separately as an alternative regime dominated by quadratic curvature terms.
